
générateurs de $\underline{J}$ au voisinage de $x$ tel que le complexe de Koszul associé soit acyclique (il en existe si on suppose l'immersion régulière); soit $g_1,\ldots,g_d$ un autre système minimum de générateurs de $\underline{J}$; il existe un voisinage de $x$ sur lequel $\underline{J}$ peut être engendré soit par les $f_i$, soit par les $g_j$, et par suite, il existe sur ce voisinage une matrice à coefficients dans $\underline{O}_X$ permettant d'exprimer les $g_j$ en fonction des $f_i$. Or les images des $f_i$ (resp. $g_j$) dans $\underline{J}/\underline{J}^2$ forment une base au point $x$, donc dans un voisinage de $x$, et par suite cette matrice est inversible dans un voisinage de $x$. Elle permet donc de définir un isomorphisme du complexe de Koszul défini par les $g_j$ sur le complexe de Koszul défini par les $f_i$, ce qui montre que le complexe de Koszul défini par les $g_j$ est acyclique.

\textbf{1.4.3.} En supposant seulement $\underline{J}$ de type fini, et $\underline{J}/\underline{J}^2$ localement libre de type fini, notons que tout système minimum de générateurs de $\underline{J}$ au voisinage d'un point $y$ de $Y$ donne une base de $\underline{J}/\underline{J}^2$ au voisinage de $y$. Réciproquement, si on a au voisinage de $y$ une base de $\underline{J}/\underline{J}^2$, on peut la relever au voisinage de $y$ en une famille de sections de $\underline{J}$, qui engendrent $\underline{J}_y$ par Nakayama, donc $\underline{J}$ dans un voisinage de $y$ puisque $\underline{J}$ est de type fini; de plus ces générateurs sont en nombre minimum, puisqu'ils donnent une base de $\underline{J}/\underline{J}^2$.

\textbf{Proposition 1.5.} (\uline{Comparer} EGA IV 19.1.5 (ii)). \uline{Soit} $i:Y\longrightarrow X$ \uline{une immersion}. \uline{Si} $i$ \uline{est régulière}, \uline{il en est de même pour toute immersion déduite de} $i$ \uline{par changement de base plat}. \uline{Pour que} $i$ \uline{soit régulière}, \uline{il faut et il suffit qu'elle le soit localement pour la topologie fidèlement plate quasi-compacte}.

La stabilité par changement de base plat est évidente.

Réciproquement, soit $i'$ l'immersion déduite de $i$ par un changement de
